% =====================================================================
%  MAC0426 / MAC5760 - Projeto 2 (Artigo sobre Topico de BD)
%  Tema: Overhead de Operacoes Criptograficas em Bancos de Dados
%        Distribuidos
%  Template: SBC (Sociedade Brasileira de Computacao) - conferencias
%
%  COMO COMPILAR (local): pdflatex -> bibtex -> pdflatex -> pdflatex
%  COMO COMPILAR (Overleaf): ver README.md. Compilador: pdfLaTeX.
%
%  ESTRUTURA: cada \section abaixo traz um comentario "% GUIA:" que
%  explica o que deve ser escrito ali e a qual item do enunciado
%  (Secao 1 do PDF do projeto) ele corresponde. Os trechos marcados
%  com "% [PREENCHER]" sao os que o grupo deve completar.
% =====================================================================

\documentclass[12pt]{article}

\usepackage{sbc-template}

\usepackage{graphicx,url}
\usepackage[utf8]{inputenc}     % arquivo salvo em UTF-8
\usepackage[T1]{fontenc}
\usepackage[brazil]{babel}      % artigo em portugues
\usepackage{booktabs}           % tabelas com regras horizontais limpas
\usepackage{amsmath}

\sloppy

% ---------------------------------------------------------------------
% TITULO E AUTORES
% ---------------------------------------------------------------------
\title{Overhead de Operações Criptográficas em\\ Bancos de Dados Distribuídos}

\author{Beatriz Viana Costa\inst{1}, Lys Katherine Park Kang\inst{1},\\
  Raphaella Brandão Jacques\inst{1}, Stephanie Maria Braga\inst{1}}

\address{Instituto de Matemática e Estatística (IME)\\
  Universidade de São Paulo (USP) -- São Paulo, SP -- Brasil
}

\begin{document}

\maketitle

% ---------------------------------------------------------------------
% RESUMO / ABSTRACT  (item 1: contextualizacao; primeira pagina; <=10 linhas)
% Rascunho inicial derivado da descricao enviada na definicao do tema.
% Refinar ao final, quando os resultados estiverem consolidados.
% ---------------------------------------------------------------------
\begin{abstract}
  As distributed systems increasingly store sensitive data, cryptographic
  techniques become essential to guarantee confidentiality and integrity.
  However, these techniques introduce performance costs that are amplified in
  distributed settings, where inter-node communication is already a natural
  bottleneck. This work investigates the impact of cryptographic operations on
  the performance of distributed databases, combining a literature review with
  practical experiments. We configure a distributed database and measure
  latency and throughput under different scenarios, with and without encryption
  (in transit and at rest), in order to quantify the introduced overhead and
  discuss the trade-offs between security and performance.
  % [PREENCHER: ajustar para refletir os resultados obtidos]
\end{abstract}

\begin{resumo}
  À medida que sistemas distribuídos passam a armazenar volumes crescentes de
  dados sensíveis, técnicas criptográficas tornam-se indispensáveis para
  garantir confidencialidade e integridade. Contudo, essas técnicas introduzem
  custos de desempenho que se tornam ainda mais relevantes em ambientes
  distribuídos, nos quais a comunicação entre nós já representa um gargalo
  natural. Este trabalho investiga o impacto de operações criptográficas sobre
  o desempenho de bancos de dados distribuídos, combinando uma revisão da
  literatura com experimentos práticos. Configuramos um banco de dados
  distribuído e medimos latência e vazão sob diferentes cenários, com e sem
  criptografia (em trânsito e em repouso), a fim de quantificar o overhead
  introduzido e discutir os trade-offs entre segurança e desempenho.
  % [PREENCHER: ajustar para refletir os resultados obtidos]
\end{resumo}


% =====================================================================
\section{Introdução}
\label{sec:introducao}
% GUIA (itens 1 e 2 do enunciado):
%  - Contextualizacao e motivacao: por que criptografia em BDs distribuidos
%    importa hoje (LGPD/GDPR, dados sensiveis, nuvem, multiparty); por que o
%    overhead e' um problema especialmente em ambientes distribuidos (a
%    comunicacao entre nos ja e' gargalo).
%  - Objetivos do estudo: declarar de forma clara o que sera investigado e as
%    perguntas de pesquisa (RQs).
%  - Fechar com a "organizacao do texto" (paragrafo descrevendo as secoes).

% [PREENCHER: contextualizacao e motivacao]
A crescente adoção de bancos de dados distribuídos para o armazenamento de
grandes volumes de dados sensíveis torna a proteção criptográfica um requisito
central de projeto~\cite{fuller2017sok}. Entretanto, garantir confidencialidade
e integridade tem um custo: operações criptográficas consomem CPU e introduzem
latência adicional, efeito agravado em sistemas distribuídos, nos quais a
comunicação entre nós já constitui um gargalo natural~\cite{gilbert2002cap}.

% [PREENCHER: objetivos e perguntas de pesquisa]
Este trabalho tem como objetivo quantificar o overhead introduzido por
operações criptográficas sobre o desempenho de bancos de dados distribuídos e
discutir o trade-off entre segurança e desempenho. Especificamente, buscamos
responder:
\begin{itemize}
  \item \textbf{PP1.} Qual o impacto da criptografia \emph{em trânsito} (TLS)
        sobre latência e vazão? % [PREENCHER/ajustar]
  \item \textbf{PP2.} Qual o impacto da criptografia \emph{em repouso} (TDE)?
  \item \textbf{PP3.} Como esse overhead varia com a carga, o número de nós e o
        perfil de leitura/escrita?
\end{itemize}

% [PREENCHER: organizacao do texto]
O restante do artigo está organizado como segue. A
Seção~\ref{sec:conceitos} apresenta os conceitos fundamentais. A
Seção~\ref{sec:relacionados} discute o estado da arte. A
Seção~\ref{sec:metodologia} detalha a metodologia experimental. A
Seção~\ref{sec:resultados} apresenta os resultados, discutidos na
Seção~\ref{sec:discussao}. A Seção~\ref{sec:conclusao} conclui o trabalho.


% =====================================================================
\section{Conceitos Fundamentais}
\label{sec:conceitos}

Para compreender o impacto de operações criptográficas em bancos de dados
distribuídos, é necessário primeiro entender como esses sistemas funcionam,
quais mecanismos de criptografia estão disponíveis e como medimos seu
desempenho. Esta seção apresenta esses três pilares.

% ---------------------------------------------------------------------
\subsection{Bancos de Dados Distribuídos}
\label{sec:conc-bdd}

Um \textbf{Banco de Dados Distribuído} (BDD) é uma coleção de bancos de dados
logicamente inter-relacionados, fisicamente dispersos por uma rede de
computadores e gerenciados por um Sistema de Gerenciamento de Banco de Dados
Distribuído (SGBDD)~\cite{elmasri2016fundamentals}. A motivação para distribuir
dados inclui maior disponibilidade, melhor desempenho por proximidade
geográfica e escalabilidade horizontal. Contudo, essa distribuição introduz
desafios que não existem em sistemas centralizados – em particular, a
necessidade de comunicação constante entre nós, que é justamente o canal sobre
o qual a criptografia em trânsito atua.

\subsubsection{Fragmentação e Replicação}

Para distribuir os dados entre os nós, um BDD utiliza duas técnicas
complementares: fragmentação e replicação~\cite{elmasri2016fundamentals}.

A \textbf{fragmentação} (também chamada de \emph{sharding} ou particionamento)
consiste em dividir uma relação em partes menores, cada uma alocada em um nó
diferente. Há três abordagens: a \emph{fragmentação horizontal} divide as
tuplas segundo algum critério de seleção (por exemplo, por região geográfica);
a \emph{fragmentação vertical} distribui os atributos da relação, como em uma operação de projeção; e a
\emph{fragmentação mista} combina ambas as estratégias. A escolha da estratégia
influencia diretamente o volume de comunicação entre nós – uma consulta que
precisa recombinar fragmentos dispersos gera tráfego de rede, e portanto mais
dados a serem cifrados quando TLS está habilitado.

A \textbf{replicação}, por sua vez, mantém cópias dos dados em múltiplos nós.
Isso melhora a disponibilidade (se um nó falha, outro responde) e o desempenho
de leitura (requisições podem ser atendidas localmente). Em contrapartida,
escritas precisam ser propagadas para todas as réplicas, gerando tráfego
adicional de sincronização. Em um cenário com TLS, cada mensagem de
sincronização entre réplicas também passa por cifragem, amplificando o overhead
proporcionalmente ao grau de replicação.

\subsubsection{Comunicação e Processamento de Consultas}

Em um BDD, o processamento de consultas difere fundamentalmente do caso
centralizado: o otimizador precisa considerar não apenas o custo de I/O e CPU,
mas também o \textbf{custo de comunicação} – isto é, o volume de dados
transferidos pela rede entre os nós~\cite{elmasri2016fundamentals}. Estratégias
como a semi-junção (\emph{semi-join}) existem justamente para reduzir esse
tráfego. Esse custo de comunicação é o ponto central deste trabalho: quando
habilitamos TLS, cada byte transmitido entre nós precisa ser cifrado na origem
e decifrado no destino, adicionando latência a cada troca de mensagem.

Além disso, para garantir a atomicidade de transações que envolvem múltiplos
nós, BDDs utilizam protocolos de coordenação como o \emph{Two-Phase Commit}
(2PC), em que um coordenador troca mensagens com todos os participantes para
decidir se a transação será confirmada ou abortada~\cite{elmasri2016fundamentals}.
Protocolos de consenso mais modernos, como
Raft~\cite{ongaro2014raft}, são utilizados em sistemas como
Spanner~\cite{corbett2012spanner}. Todos esses protocolos dependem de múltiplas
rodadas de comunicação, o que os torna sensíveis ao overhead de criptografia em
trânsito.

\subsubsection{Teorema CAP}

O teorema CAP~\cite{gilbert2002cap} formaliza um trade-off fundamental: em um
sistema distribuído sujeito a partições de rede, é impossível garantir
simultaneamente \emph{Consistência} (todos os nós veem os mesmos dados),
\emph{Disponibilidade} (toda requisição recebe resposta) e \emph{Tolerância a
Partições}. Como partições são inevitáveis na prática, sistemas precisam optar
entre priorizar consistência (modelo CP, como Spanner) ou disponibilidade
(modelo AP, como Cassandra). Essa escolha arquitetural determina a quantidade
e o padrão de comunicação entre nós, influenciando diretamente o quanto o
overhead do TLS se manifesta no desempenho global do sistema.

% ---------------------------------------------------------------------
\subsection{Técnicas Criptográficas em Bancos de Dados}
\label{sec:conc-cripto}

A proteção de dados em um banco de dados pode ocorrer em três momentos
distintos, de acordo com o estado em que os dados se
encontram~\cite{nist2020sp80053r5}.

\subsubsection{Criptografia em Repouso (\emph{at rest})}

Quando os dados estão armazenados em disco, a criptografia em repouso protege
contra acessos não autorizados ao meio físico (roubo de disco, acesso indevido
a arquivos do sistema operacional). A abordagem mais comum é a
\emph{Transparent Data Encryption} (TDE), que cifra as páginas de dados e os
arquivos de log de forma transparente para a aplicação – ou seja, o banco
realiza a cifragem e decifragem automaticamente nas operações de I/O, sem
alterar as consultas~\cite{nist2020sp80053r5}.

O algoritmo predominante é o AES (\emph{Advanced Encryption Standard}), um
cifrador de bloco simétrico padronizado pelo NIST~\cite{nist2001aes}, com
chaves de 128, 192 ou 256 bits. A granularidade pode variar: criptografia de
todo o disco, por \emph{tablespace}, por tabela ou mesmo por coluna individual.
O overhead da criptografia em repouso incide sobre cada operação de leitura e
escrita em disco, pois os dados precisam ser decifrados ao serem carregados para
a memória e cifrados novamente ao serem persistidos.

\subsubsection{Criptografia em Trânsito (\emph{in transit})}

A criptografia em trânsito protege os dados enquanto trafegam pela rede –
seja entre um cliente e o servidor, seja entre os nós internos do cluster
distribuído. O protocolo padrão é o TLS (\emph{Transport Layer
Security})~\cite{rescorla2018tls}, atualmente na versão 1.3.

O custo de desempenho do TLS pode ser decomposto em dois componentes:

\begin{enumerate}
  \item \textbf{Handshake}: no estabelecimento da conexão, cliente e servidor
        negociam parâmetros criptográficos utilizando criptografia assimétrica
        (RSA ou curvas elípticas). É uma operação computacionalmente cara, mas
        realizada apenas uma vez por conexão (ou por reconexão).
  \item \textbf{Cifragem simétrica contínua}: após o handshake, todos os dados
        são cifrados com algoritmos simétricos como AES-GCM ou
        ChaCha20-Poly1305. O custo por pacote é baixo individualmente, mas em
        cenários de alta vazão (muitas operações por segundo ou grandes volumes
        de dados) o impacto acumula.
\end{enumerate}

Em um banco de dados distribuído, o TLS pode ser habilitado tanto na
comunicação cliente--servidor quanto na comunicação interna entre os nós do
cluster (\emph{intra-cluster}). Neste segundo caso, o impacto tende a ser mais
expressivo, pois protocolos de replicação e consenso geram tráfego constante
entre os nós.

\subsubsection{Criptografia em Uso (\emph{in use})}

Uma linha de pesquisa mais recente busca permitir computações diretamente sobre
dados cifrados, eliminando a necessidade de decifrá-los. As abordagens
incluem: criptografia determinística e que preserva ordem
(OPE)~\cite{boldyreva2009ope}, que viabilizam consultas de igualdade e
intervalos sobre dados cifrados, como explorado pelo sistema
CryptDB~\cite{popa2011cryptdb}; criptografia totalmente homomórfica
(FHE)~\cite{gentry2009fhe}, que permite computações arbitrárias mas com
overhead ainda proibitivo para cargas transacionais; e soluções baseadas em
enclaves seguros (\emph{trusted execution environments}), como o Always
Encrypted~\cite{antonopoulos2020alwaysencrypted}. Embora não sejam o foco dos
nossos experimentos, essas técnicas representam o estado da arte em proteção de
dados e servem como horizonte para trabalhos futuros.

\subsubsection{Integridade}

Complementando a confidencialidade, mecanismos de integridade garantem que os
dados não foram corrompidos ou adulterados. Técnicas como \emph{hashing}
criptográfico (SHA-256), \emph{HMAC} (Hash-based Message Authentication Code) e
árvores de Merkle permitem detectar alterações tanto em dados armazenados quanto
em dados em trânsito~\cite{nist2020sp80053r5}. No contexto do TLS, a
integridade é garantida implicitamente pelos modos de operação AEAD
(\emph{Authenticated Encryption with Associated Data}), como o AES-GCM, que
combina cifragem e autenticação em uma única operação.

% ---------------------------------------------------------------------
\subsection{Métricas de Desempenho}
\label{sec:conc-metricas}

Para avaliar o impacto das operações criptográficas de forma rigorosa, é
necessário definir as métricas que serão
observadas~\cite{cooper2010ycsb}:

\begin{itemize}
  \item \textbf{Latência}: o tempo entre o envio de uma requisição e o
        recebimento da resposta correspondente. Como a distribuição de latências
        tende a ser assimétrica (com uma ``cauda longa''), reportamos não apenas
        a média, mas também a mediana (p50) e os percentis p95 e p99, que
        revelam o comportamento nos piores casos.
  \item \textbf{Vazão (\emph{throughput})}: o número de operações completadas
        por unidade de tempo (operações/segundo). Enquanto a latência captura a
        experiência individual de cada requisição, a vazão reflete a capacidade
        agregada do sistema sob carga.
  \item \textbf{Utilização de CPU}: a fração do processamento dedicada a
        operações criptográficas. Permite distinguir se o overhead observado
        decorre de cifragem ou de outros fatores (I/O, contenção de rede).
  \item \textbf{Overhead de rede}: o aumento no volume de bytes transmitidos
        devido aos cabeçalhos e metadados introduzidos pelo TLS (registros TLS,
        MACs, padding).
\end{itemize}

Para quantificar o impacto relativo da criptografia, definimos o
\textbf{overhead} em relação ao cenário \emph{baseline} (sem criptografia):
\begin{equation}
  \text{overhead}(\%) = \frac{M_{\text{cripto}} - M_{\text{baseline}}}{M_{\text{baseline}}} \times 100
\end{equation}
onde $M$ representa o valor de uma métrica (latência ou vazão) no respectivo
cenário. Um overhead positivo em latência indica degradação; um overhead
negativo em vazão indica queda na capacidade do sistema.


% =====================================================================
\section{Trabalhos Relacionados}
\label{sec:relacionados}
% GUIA (item 5 - "estado da arte de trabalhos teoricos e, se pertinente,
% ferramentas de software"):
%  - Trabalhos teoricos/cientificos: surveys e estudos de overhead de cripto
%    em BDs (ex.: \cite{fuller2017sok}, \cite{popa2011cryptdb},
%    \cite{antonopoulos2020alwaysencrypted}).
%  - Ferramentas de software (priorizar SOFTWARE LIVRE, conforme enunciado):
%    SGBDs distribuidos (PostgreSQL/CockroachDB/YugabyteDB/MongoDB/Cassandra),
%    recursos de cripto de cada um, e ferramentas de benchmark
%    (YCSB \cite{cooper2010ycsb}, sysbench, BenchBase/TPC-C).
% [PREENCHER]: situar o nosso estudo em relacao a esses trabalhos (lacuna).


% =====================================================================
\section{Metodologia}
\label{sec:metodologia}
% GUIA (item 3 - "metodos empregados"; o metodo aqui e' EXPERIMENTO):
%  Descrever com detalhe suficiente para REPRODUCAO (criterio do enunciado).

\subsection{Ambiente Experimental}
\label{sec:met-ambiente}
% [PREENCHER]: hardware/VM, SO, versao do SGBD, topologia (ex.: cluster de 3
% nos via Docker Compose), versoes de bibliotecas TLS. SGBD a definir.

\subsection{Carga de Trabalho}
\label{sec:met-carga}
% [PREENCHER]: ferramenta de benchmark (ex.: YCSB \cite{cooper2010ycsb} ou
% sysbench), workloads (proporcao leitura/escrita), tamanho do dataset,
% tamanho dos registros, numero de threads/clientes.

\subsection{Cenários Avaliados}
\label{sec:met-cenarios}
% [PREENCHER]: descrever os cenarios. A Tabela~\ref{tab:cenarios} resume.

% --- Exemplo de tabela (estilo SBC: legenda ANTES da tabela) ------------
\begin{table}[ht]
\centering
\caption{Cenários de criptografia avaliados.}
\label{tab:cenarios}
\begin{tabular}{@{}lll@{}}
\toprule
\textbf{Cenário} & \textbf{Em trânsito (TLS)} & \textbf{Em repouso (TDE)} \\
\midrule
C0 (baseline)    & não & não \\
C1               & sim & não \\
C2               & não & sim \\
C3               & sim & sim \\
\bottomrule
\end{tabular}
\end{table}

\subsection{Protocolo de Medição}
\label{sec:met-protocolo}
% [PREENCHER]: numero de repeticoes, tempo de warm-up, descarte de outliers,
% como a media e o desvio-padrao sao calculados, como o overhead relativo
% e' computado em relacao ao C0. Garantir reprodutibilidade.


% =====================================================================
\section{Experimentos e Resultados}
\label{sec:resultados}
% GUIA (item 6 - "descricao do estudo realizado e apresentacao dos resultados"):
%  - Apresentar os dados de forma clara (graficos + tabelas).
%  - Reportar media +/- desvio; nao usar mais casas decimais que a precisao
%    justifica (recomendacao do template SBC).

% --- Exemplo de figura (gerar o grafico e salvar em figuras/) -----------
% \begin{figure}[ht]
% \centering
% \includegraphics[width=.7\textwidth]{figuras/vazao_por_cenario.pdf}
% \caption{Vazão (ops/s) por cenário de criptografia.}
% \label{fig:vazao}
% \end{figure}

% [PREENCHER]: vazao por cenario; latencia (p50/p95/p99) por cenario;
% overhead relativo (%) em relacao ao C0; variacao com numero de nos/carga.


% =====================================================================
\section{Discussão}
\label{sec:discussao}
% GUIA (item 7 - "discussao sobre os resultados"):
% [PREENCHER]: interpretar o trade-off seguranca x desempenho; quando o
% overhead é aceitavel; qual componente domina (handshake vs. cifra vs. I/O);
% ameacas a validade (internas/externas) e limitacoes do estudo.


% =====================================================================
\section{Conclusão e Trabalhos Futuros}
\label{sec:conclusao}
% [PREENCHER]: sintetizar as respostas as perguntas de pesquisa e apontar
% extensoes (ex.: criptografia em uso / queryable encryption, enclaves,
% outros SGBDs, cargas analiticas (OLAP)).


% =====================================================================
% REFERENCIAS  (item: qualidade das referencias e' criterio de avaliacao)
% Estilo SBC (sbc.bst). Adicionar/editar entradas em referencias.bib.
% =====================================================================
\bibliographystyle{sbc}
\bibliography{referencias}

\end{document}
